% Telos shared preamble.
%
% Every Telos publication is designed to be printed on a home laser printer,
% one sheet at a time, and carried into a boat or a garage. Two constraints
% follow from that and govern everything below:
%
%   1. Pure monochrome. No value is ever a hue. Emphasis is carried by rule
%      weight, fill, and type, never by color, so a grayscale or black-only
%      printer loses nothing.
%   2. Card density. A "one-pager" is a hard contract: the leaf must fit on a
%      single side of US Letter. The layout is therefore tight by default and
%      the environments below are sized to leave room for a plate.

\documentclass[10pt,letterpaper]{article}

\usepackage[margin=0.55in,includefoot,footskip=0.24in]{geometry}
\usepackage[T1]{fontenc}
\usepackage[utf8]{inputenc}
\usepackage{lmodern}
\usepackage{microtype}
\usepackage{array}
\usepackage{booktabs}
\usepackage{longtable}
\usepackage{tabularx}
\usepackage{enumitem}
\usepackage{needspace}
\usepackage{multicol}
\usepackage{ragged2e}
\usepackage{graphicx}
\usepackage{wrapfig}
\usepackage{xcolor}
\usepackage{fancyhdr}
\usepackage{titlesec}
\usepackage{hyperref}
\usepackage[most]{tcolorbox}
\usepackage{tikz}
\usetikzlibrary{arrows.meta,calc,positioning,shapes.geometric,decorations.pathmorphing,patterns,fadings}

% Semantic names are kept so a document reads intelligibly, but every value is
% black, white, or a neutral gray that survives a black-only print driver.
\definecolor{telosink}{gray}{0}
\definecolor{telospaper}{gray}{1}
\definecolor{telosrule}{gray}{0.35}
\definecolor{telosfill}{gray}{0.90}
\definecolor{telosdeep}{gray}{0.72}
\definecolor{telosmute}{gray}{0.45}

\hypersetup{hidelinks,pdfborder={0 0 0}}

% Keep an unchanged source byte-reproducible across rebuilds: no build-clock
% creation date and no random trailer ID.
\pdfinfoomitdate=1
\pdftrailerid{}

\setlength{\parindent}{0pt}
\setlength{\parskip}{0.42em}
\setlength{\tabcolsep}{4pt}
\setlength{\columnsep}{1.1em}
\renewcommand{\arraystretch}{1.18}
\setlist[itemize]{leftmargin=1.15em,itemsep=0.12em,topsep=0.2em,parsep=0pt}
\setlist[enumerate]{leftmargin=1.5em,itemsep=0.18em,topsep=0.2em,parsep=0pt}
\setlist[description]{leftmargin=0pt,itemsep=0.16em,topsep=0.2em,parsep=0pt,
  font=\normalfont\bfseries}

% ---------------------------------------------------------------------------
% Document identity. Each leaf declares these before \begin{document} so the
% running head, the colophon, and the PDF metadata all agree.
% ---------------------------------------------------------------------------
\newcommand{\TelosProject}{Telos}
\newcommand{\TelosDocTitle}{}
\newcommand{\TelosDocSubtitle}{}
\newcommand{\TelosRevision}{}

\newcommand{\telosmeta}[4]{%
  \renewcommand{\TelosProject}{#1}%
  \renewcommand{\TelosDocTitle}{#2}%
  \renewcommand{\TelosDocSubtitle}{#3}%
  \renewcommand{\TelosRevision}{#4}%
  \hypersetup{pdftitle={#2},pdfsubject={#3; version #4},
    pdfkeywords={Telos, version #4},pdfauthor={Telos}}%
}

\pagestyle{fancy}
\fancyhf{}
\renewcommand{\headrulewidth}{0pt}
\renewcommand{\footrulewidth}{0.4pt}
\fancyfoot[L]{\footnotesize\textsc{\TelosProject}}
\fancyfoot[C]{\footnotesize\TelosDocTitle}
\fancyfoot[R]{\footnotesize\thepage}

% The first page of a one-pager carries the masthead instead of a running foot,
% so its footer is the compact provenance line.
\fancypagestyle{telosfirst}{%
  \fancyhf{}%
  \renewcommand{\headrulewidth}{0pt}%
  \renewcommand{\footrulewidth}{0.4pt}%
  \fancyfoot[L]{\footnotesize\textsc{\TelosProject}}%
  \fancyfoot[C]{\footnotesize\TelosRevision}%
  \fancyfoot[R]{\footnotesize\thepage}%
}

% ---------------------------------------------------------------------------
% Masthead. A heavy rule, the title, a light subtitle, a thin rule. The
% optional argument is a right-aligned tag (a species' scientific name, a
% lesson number) that sits on the title's baseline.
% ---------------------------------------------------------------------------
\newcommand{\telosmasthead}[1][]{%
  \thispagestyle{telosfirst}%
  \begingroup
  \setlength{\parskip}{0pt}%
  \noindent\rule{\linewidth}{1.6pt}\par
  \vspace{0.30em}
  \noindent
  \begin{minipage}[b]{0.66\linewidth}
    \raggedright{\LARGE\bfseries\TelosDocTitle}
  \end{minipage}%
  \hfill
  \begin{minipage}[b]{0.32\linewidth}
    \raggedleft{\small\itshape #1}
  \end{minipage}\par
  \vspace{0.18em}
  \noindent{\small\TelosDocSubtitle}\par
  \vspace{0.28em}
  \noindent\rule{\linewidth}{0.4pt}\par
  \endgroup
  \vspace{0.35em}
}

% ---------------------------------------------------------------------------
% Headings. Compact, small-caps, rule-anchored — a card has no room for the
% article class's default vertical air.
% ---------------------------------------------------------------------------
\titleformat{\section}{\normalfont\large\bfseries}{}{0pt}{}[\vspace{-0.55em}\rule{\linewidth}{0.4pt}]
\titlespacing*{\section}{0pt}{0.85em}{0.35em}
\titleformat{\subsection}{\normalfont\normalsize\bfseries}{}{0pt}{}
\titlespacing*{\subsection}{0pt}{0.6em}{0.2em}
\titleformat{\subsubsection}{\normalfont\small\bfseries\scshape}{}{0pt}{}
\titlespacing*{\subsubsection}{0pt}{0.45em}{0.12em}

% A short run-in heading for card interiors that must not consume a full line.
\newcommand{\telosrubric}[1]{\textbf{#1}\enspace}

% Cross-reference helper. Every Telos project is a set of loose printed sheets,
% so a reference names the sheet rather than a page number.
\newcommand{\sheetref}[1]{\textit{see the \textbf{#1} sheet}}

% Which decision, ADR or sheet governs the paragraph you are reading. A reader
% who disagrees with an instruction should be able to find the reasoning behind
% it rather than having to take it on trust.
\newcommand{\governedby}[1]{{\footnotesize\textsc{governed by #1}}}

% ---------------------------------------------------------------------------
% Quantities. siunitx is not in this TeX Live split, and the guides need only a
% fixed thin space and upright units, so define that directly rather than take
% a package dependency the documented install would not satisfy.
%   \qty{9}{kV}  ->  9 kV        \unitof{kV}  ->  kV
% ---------------------------------------------------------------------------
\newcommand{\unitof}[1]{\ensuremath{\mathrm{#1}}}
\newcommand{\qty}[2]{#1\,\unitof{#2}}
\newcommand{\numrange}[2]{#1\,--\,#2}
\newcommand{\qtyrange}[3]{#1\,--\,#2\,\unitof{#3}}
% Inches and feet appear constantly in the fishing guide; keep one spelling.
\newcommand{\inch}[1]{#1\,in}
\newcommand{\feet}[1]{#1\,ft}

% ---------------------------------------------------------------------------
% Cards and boxes.
% ---------------------------------------------------------------------------

% Titled card with a hairline frame. The default card for facts a reader looks
% up rather than reads through.
\newtcolorbox{teloscard}[1]{
  enhanced, breakable,
  colback=telospaper, colframe=telosink, colbacktitle=telospaper,
  coltitle=telosink,
  boxrule=0.5pt, sharp corners,
  left=6pt, right=6pt, top=4pt, bottom=4pt,
  toptitle=3pt, bottomtitle=3pt, titlerule=0.4pt,
  before skip=6pt, after skip=6pt,
  title={#1}, fonttitle=\small\bfseries
}

% Untitled left-ruled block: hierarchy without a redundant label.
\newtcolorbox{telosframe}{
  enhanced, breakable,
  colback=telospaper, frame hidden, boxrule=0pt, sharp corners,
  borderline west={1.3pt}{0pt}{telosink},
  left=8pt, right=2pt, top=3pt, bottom=3pt,
  before skip=6pt, after skip=6pt
}

% Filled emphasis panel — the "read this first" block. The 10% fill stays
% legible under a cheap laser toner curve and still reads as distinct.
\newtcolorbox{telospanel}[1]{
  enhanced, breakable,
  colback=telosfill, colframe=telosink, colbacktitle=telosfill,
  coltitle=telosink,
  boxrule=0.5pt, sharp corners,
  left=6pt, right=6pt, top=4pt, bottom=4pt,
  toptitle=3pt, bottomtitle=3pt, titlerule=0.4pt,
  before skip=6pt, after skip=6pt,
  title={#1}, fonttitle=\small\bfseries
}

% Hazard block. Deliberately the heaviest object on any page: a 1.4pt frame
% plus a solid title bar. Nothing else in Telos is allowed to look like this,
% so a reader skimming a lesson cannot mistake it for ordinary emphasis.
\newtcolorbox{teloshazard}[1]{
  enhanced, breakable,
  colback=telospaper, colframe=telosink,
  colbacktitle=telosink, coltitle=telospaper,
  boxrule=1.4pt, sharp corners,
  left=6pt, right=6pt, top=5pt, bottom=5pt,
  toptitle=3pt, bottomtitle=3pt,
  before skip=8pt, after skip=8pt,
  title={\MakeUppercase{#1}}, fonttitle=\small\bfseries
}

% A stop-and-check gate. Used before anything destructive, and before anything
% that would be expensive to discover later. Shared by every project that has
% a procedure someone could get hurt following carelessly.
\newtcolorbox{gate}[1]{
  enhanced, breakable,
  colback=telospaper, colframe=telosink,
  boxrule=1.0pt, sharp corners,
  left=6pt, right=6pt, top=4pt, bottom=4pt,
  toptitle=3pt, bottomtitle=3pt, titlerule=0.4pt,
  before skip=7pt, after skip=7pt,
  colbacktitle=telospaper, coltitle=telosink,
  title={\MakeUppercase{#1}}, fonttitle=\small\bfseries
}

% ---------------------------------------------------------------------------
% Rating dots. A five-step scale that prints identically in black-only: filled
% versus open circles, never a shaded bar.
% ---------------------------------------------------------------------------
\newcommand{\telosdot}{\tikz[baseline=-0.42ex]\fill[telosink] (0,0) circle (0.28ex);}
\newcommand{\telosopendot}{\tikz[baseline=-0.42ex]\draw[telosink,line width=0.28pt] (0,0) circle (0.28ex);}
\makeatletter
\newcounter{telos@rate}
\newcommand{\telosrating}[1]{%
  \begingroup
  \setcounter{telos@rate}{0}%
  \@whilenum\value{telos@rate}<#1\do{\telosdot\kern0.22em\stepcounter{telos@rate}}%
  \@whilenum\value{telos@rate}<5\do{\telosopendot\kern0.22em\stepcounter{telos@rate}}%
  \endgroup
}
\makeatother

% ---------------------------------------------------------------------------
% Tables.
%
% tabularx reads its body by scanning for a literal \end{tabularx}, so it can
% never be wrapped in a \newenvironment. Documents therefore open tabularx and
% longtable directly and take their house style from these column types and the
% booktabs rules. The one exception is the fact strip below, which is built on
% plain tabular precisely so it *can* be an environment.
% ---------------------------------------------------------------------------
\newcolumntype{L}{>{\raggedright\arraybackslash}X}
\newcolumntype{R}{>{\raggedleft\arraybackslash}X}
\newcolumntype{C}{>{\centering\arraybackslash}X}
\newcolumntype{B}{>{\bfseries\raggedright\arraybackslash}X}
\newcolumntype{N}{>{\raggedright\arraybackslash}p{0.16\linewidth}}

% Key/value fact strip. Two flush columns of short lookups; the label column is
% a fixed fraction so a stack of cards aligns across the whole guide.
\newenvironment{telosfacts}[1][0.24]
  {\begingroup\small
   \setlength{\parskip}{0pt}%
   \renewcommand{\arraystretch}{1.25}%
   \begin{tabular}{@{}>{\bfseries\raggedright\arraybackslash}p{\dimexpr#1\linewidth\relax}@{\hspace{0.7em}}>{\raggedright\arraybackslash}p{\dimexpr\linewidth-#1\linewidth-0.7em\relax}@{}}}
  {\end{tabular}\par\endgroup}
\newcommand{\telosfact}[2]{#1 & #2\\}

% ---------------------------------------------------------------------------
% Seasonal / diurnal activity bar.
%
% \telosseasonbar{4,3,2,1,1,2,3,3,4,5,4,3} draws a twelve-cell month strip
% where each value 0-5 is a fill height. Height, not gray value, carries the
% signal, so it is exact under any toner curve; the cell is outlined either way
% so an empty month is still visibly a month.
% ---------------------------------------------------------------------------
\newcommand{\telosseasonlabels}{J,F,M,A,M,J,J,A,S,O,N,D}
\newcommand{\telosseasonbar}[1]{%
  \begingroup
  \begin{tikzpicture}[x=1.25em,y=2.1ex,baseline=0pt]
    \foreach \v [count=\i from 0] in {#1} {
      \draw[telosrule,line width=0.3pt] (\i,0) rectangle (\i+0.86,5);
      \ifnum\v>0
        \fill[telosink] (\i,0) rectangle (\i+0.86,\v);
      \fi
    }
    \foreach \m [count=\i from 0] in \telosseasonlabels {
      \node[font=\tiny,anchor=north,inner sep=1.2pt] at (\i+0.43,-0.1) {\m};
    }
  \end{tikzpicture}%
  \endgroup
}

% Four-cell day-part strip: dawn, midday, dusk, night.
\newcommand{\telosdaylabels}{Dawn,Mid,Dusk,Night}
\newcommand{\telosdaybar}[1]{%
  \begingroup
  \begin{tikzpicture}[x=2.9em,y=2.1ex,baseline=0pt]
    \foreach \v [count=\i from 0] in {#1} {
      \draw[telosrule,line width=0.3pt] (\i,0) rectangle (\i+0.92,5);
      \ifnum\v>0
        \fill[telosink] (\i,0) rectangle (\i+0.92,\v);
      \fi
    }
    \foreach \m [count=\i from 0] in \telosdaylabels {
      \node[font=\tiny,anchor=north,inner sep=1.2pt] at (\i+0.46,-0.1) {\m};
    }
  \end{tikzpicture}%
  \endgroup
}

% A bar needs vertical room for its labels when it sits in a fact strip.
\newcommand{\telosbarrow}[1]{\rule{0pt}{4.2ex}#1\rule[-1.6ex]{0pt}{1.6ex}}

% ---------------------------------------------------------------------------
% Plates. A species or apparatus illustration with a caption rule beneath it.
% \telosplate{width fraction}{file}{caption}
% ---------------------------------------------------------------------------
\newcommand{\telosplate}[3]{%
  \begingroup
  \centering
  \includegraphics[width=#1\linewidth]{#2}\par
  \vspace{0.15em}
  \begin{minipage}{#1\linewidth}
    \rule{\linewidth}{0.4pt}\par
    \vspace{0.1em}
    {\footnotesize #3\par}
  \end{minipage}\par
  \endgroup
  \vspace{0.3em}
}

% TikZ drawing conventions shared by every rig, knot, and circuit diagram.
\tikzset{
  telos line/.style={draw=telosink,line width=0.7pt},
  telos hair/.style={draw=telosink,line width=0.35pt},
  telos heavy/.style={draw=telosink,line width=1.2pt},
  telos node/.style={draw=telosink,fill=telospaper,align=center,inner sep=4pt,font=\small},
  telos emphasis/.style={draw=telosink,line width=1.2pt,fill=telospaper,align=center,
    inner sep=4pt,font=\small\bfseries},
  telos label/.style={font=\scriptsize,inner sep=1.5pt},
  telos arrow/.style={-{Stealth[length=1.8mm]},draw=telosink,line width=0.6pt},
  telos dim/.style={|<->|,draw=telosink,line width=0.35pt},
  telos water/.style={draw=telosrule,line width=0.4pt,decorate,
    decoration={snake,amplitude=0.4mm,segment length=3mm}},
}

% ---------------------------------------------------------------------------
% Lesson furniture
% ---------------------------------------------------------------------------

% \lessonhead{number}{time}{who does what}
\newcommand{\lessonhead}[3]{%
  \par\noindent
  \begin{minipage}{\linewidth}
    \rule{\linewidth}{0.4pt}\par
    \vspace{0.3em}
    {\small\textbf{Lesson #1}\quad$\cdot$\quad #2\quad$\cdot$\quad #3\par}
    \vspace{0.25em}
    \rule{\linewidth}{0.4pt}
  \end{minipage}\par
  \vspace{0.5em}}

% The materials list. Two columns of short lines; a lesson that needs a long
% shopping list is a lesson that has not been thought through.
\newenvironment{materials}
  {\begin{telospanel}{What you need}\begin{multicols}{2}
   \begin{itemize}[leftmargin=1em,itemsep=0.05em,topsep=0pt]}
  {\end{itemize}\end{multicols}\end{telospanel}}

% The four rubrics every lesson carries, in the same order every time:
% what you are about to see, what to do, what it means, and what to ask.
\newcommand{\lessonsection}[1]{\section*{#1}\vspace{-0.35em}\rule{\linewidth}{0.4pt}\vspace{0.2em}}

% A prediction box. The point of the curriculum is that they commit to an
% answer before the demonstration, not after.
\newtcolorbox{predict}{
  enhanced, breakable,
  colback=telospaper, colframe=telosink,
  boxrule=0.5pt, sharp corners,
  left=6pt, right=6pt, top=4pt, bottom=4pt,
  before skip=6pt, after skip=6pt,
  borderline west={2.2pt}{0pt}{telosink},
  title={\small\bfseries Predict first}, fonttitle=\small\bfseries,
  colbacktitle=telospaper, coltitle=telosink,
  toptitle=3pt, bottomtitle=3pt, titlerule=0.4pt,
}

% ---------------------------------------------------------------------------
% Colophon. Rendered once, at the end of the last page, in the recovered footer
% area so it never creates a page of its own.
% ---------------------------------------------------------------------------
\newif\ifTelosColophonRendered
\newcommand{\TelosColophonExtra}{}
\newcommand{\TelosColophon}{%
  \ifTelosColophonRendered\else
  \global\TelosColophonRenderedtrue
  \par
  \enlargethispage{0.5in}%
  \vspace{0.3em}%
  \begingroup
  \setlength{\parskip}{0pt}%
  \begin{minipage}{\linewidth}
  \rule{\linewidth}{0.35pt}\par
  \vspace{0.3em}%
  \fontsize{7}{8.2}\selectfont
  \textbf{Telos.} \TelosProject{} — \TelosDocTitle. Revised \TelosRevision.
  \TelosColophonExtra{}
  Project-created text and diagrams are licensed \href{https://creativecommons.org/licenses/by/4.0/}{CC BY 4.0}.
  Incorporated public-domain artwork remains public domain and is credited on
  the plate it appears with. See \texttt{LICENSE} and \texttt{CREDITS.md} in the
  source. Nothing here is professional advice; verify current regulations and
  follow the stated safety procedure.
  \end{minipage}
  \endgroup
  \fi
}
\AtEndDocument{\TelosColophon}

% ChatGPT edition map primitives. Original monochrome drawings; no Claude map
% geometry or generated output is imported.
\tikzset{
  cgshore/.style={draw=black,line width=1pt,fill=white},
  cgten/.style={draw=black,line width=.5pt},
  cgtwenty/.style={draw=black,line width=.42pt},
  cgforty/.style={draw=black,line width=.35pt,dash pattern=on 4pt off 1.5pt},
  cgsixty/.style={draw=black,line width=.32pt,dash pattern=on 2pt off 1.2pt},
  cgspot/.style={circle,fill=black,text=white,minimum size=3.3mm,inner sep=0pt,
    font=\fontsize{6}{6}\selectfont\bfseries},
  cgdepth/.style={fill=white,inner sep=.6pt,font=\fontsize{5.5}{6}\selectfont},
  cgnote/.style={fill=white,inner sep=1pt,font=\fontsize{6}{7}\selectfont}
}
\newcommand{\cgspot}[3]{\node[cgspot] at (#1,#2) {#3};}
\newcommand{\cgnorth}[2]{%
  \draw[-{Stealth[length=1.8mm]},line width=.5pt] (#1,#2) -- ++(0,.65);
  \node[font=\scriptsize\bfseries,anchor=south] at (#1,#2+.65) {N};}
\newcommand{\cgscale}[4]{%
  \draw[line width=.5pt] (#1,#2) -- ++(#3,0);
  \draw[line width=.5pt] (#1,#2-.07) -- ++(0,.14);
  \draw[line width=.5pt] (#1+#3,#2-.07) -- ++(0,.14);
  \node[font=\fontsize{5.5}{6}\selectfont,anchor=north] at (#1+#3/2,#2-.08) {#4};}

\telosmeta{Lake Country Fishing}{Pine Lake}
  {WBIC 779200 --- independent field sheet}{2026-07-26}

\begin{document}
\telosmasthead[Waukesha County, Wisconsin]

Pine is a large, deep spring lake with an overwhelmingly gravel bottom. Its
1955 survey shows sharp basin edges, a mid-lake island and long pieces of hard
structure. Use this sheet to form a search plan, then let lawful electronics,
visible vegetation, weather and bites revise it.

\begin{telosfacts}[0.19]
\telosfact{Identity}{WBIC 779200; 711 acres}
\telosfact{Depth}{85 ft maximum; 39 ft mean}
\telosfact{Bottom}{90\% gravel; 5\% sand; 5\% muck}
\telosfact{Type}{Spring lake; DNR classifies it oligotrophic}
\telosfact{Current DNR fish record}{Panfish, largemouth bass, smallmouth bass,
 northern pike and walleye}
\telosfact{Access}{DNR lists one landing. Verify the launch sign, hours, fees
 and all local ordinances on arrival}
\end{telosfacts}

\begin{center}
\begin{tikzpicture}[x=.68cm,y=.68cm]
% Independent generalization of the 1955 DNR survey, checked against 2024 NAIP.
\path[cgshore]
 (4.4,13.7)..controls(4.0,13.0)and(4.1,12.3)..(3.8,11.7)
 ..controls(3.5,11.0)and(4.0,10.4)..(3.2,9.7)
 ..controls(2.3,9.1)and(1.2,9.3)..(1.0,8.5)
 ..controls(.8,7.8)and(1.6,7.4)..(1.2,6.7)
 ..controls(.8,6.0)and(.9,5.1)..(1.6,4.7)
 ..controls(2.0,4.4)and(1.4,3.7)..(1.7,3.0)
 ..controls(2.1,2.2)and(2.7,1.6)..(3.0,.7)
 ..controls(3.4,.1)and(4.1,.5)..(4.3,1.1)
 ..controls(4.6,1.7)and(5.1,1.5)..(5.6,2.0)
 ..controls(6.2,2.5)and(6.0,3.2)..(6.6,3.7)
 ..controls(7.0,4.1)and(6.5,4.7)..(6.9,5.3)
 ..controls(7.5,6.2)and(7.1,7.1)..(7.5,8.0)
 ..controls(7.8,8.8)and(7.3,9.5)..(7.6,10.2)
 ..controls(7.9,10.9)and(7.2,11.3)..(6.8,11.1)
 ..controls(6.1,10.8)and(5.7,11.6)..(5.1,11.7)
 ..controls(4.7,11.9)and(5.0,12.8)..(4.4,13.7)--cycle;
\draw[cgten]
 (4.35,13.1)..controls(4.2,12.1)and(4.0,11.3)..(3.7,10.4)
 ..controls(3.2,9.7)and(1.6,9.3)..(1.5,8.5)
 ..controls(1.4,7.7)and(2.0,7.1)..(1.6,6.2)
 ..controls(1.2,5.1)and(2.0,4.5)..(2.0,3.6)
 ..controls(2.1,2.7)and(3.0,2.0)..(3.2,1.1)
 ..controls(3.8,.8)and(4.1,1.6)..(4.7,1.9)
 ..controls(5.6,2.3)and(5.4,3.3)..(6.2,3.9)
 ..controls(6.8,4.5)and(6.2,5.4)..(6.8,6.4)
 ..controls(7.2,7.4)and(6.8,8.7)..(7.1,9.7)
 ..controls(7.3,10.5)and(6.7,10.7)..(6.1,10.6)
 ..controls(5.2,10.4)and(5.0,11.4)..(4.6,11.2)
 ..controls(4.1,10.9)and(4.7,10.0)..(4.0,9.7)
 ..controls(3.6,9.3)and(4.1,8.7)..(4.0,8.0)
 ..controls(3.9,7.3)and(4.1,6.3)..(4.2,5.5)
 ..controls(4.4,4.6)and(4.6,3.5)..(4.1,2.7);
\draw[cgtwenty]
 (4.2,11.0)..controls(3.2,10.2)and(2.2,9.1)..(2.0,8.2)
 ..controls(1.8,7.0)and(2.4,6.4)..(2.3,5.3)
 ..controls(2.1,4.3)and(2.7,3.2)..(3.2,2.2)
 ..controls(3.8,1.6)and(4.7,2.5)..(5.0,3.3)
 ..controls(5.3,4.1)and(5.1,4.9)..(5.6,5.8)
 ..controls(6.3,7.0)and(6.0,8.3)..(6.4,9.2)
 ..controls(6.6,9.8)and(5.9,10.1)..(5.3,9.9)
 ..controls(4.8,9.7)and(5.1,8.8)..(4.6,8.3)
 ..controls(4.1,7.8)and(4.7,6.6)..(4.6,5.7)
 ..controls(4.5,4.9)and(4.2,4.1)..(4.0,3.2);
\draw[cgforty]
 (3.9,10.1)..controls(2.8,9.0)and(2.4,8.1)..(2.7,7.1)
 ..controls(3.0,6.0)and(2.6,5.0)..(3.0,4.0)
 ..controls(3.4,3.0)and(4.4,3.1)..(4.7,4.0)
 ..controls(5.0,5.0)and(5.2,5.9)..(5.5,6.9)
 ..controls(5.8,8.0)and(5.3,9.1)..(4.8,9.4)
 ..controls(4.4,9.6)and(4.2,9.9)..(3.9,10.1)--cycle;
\draw[cgsixty]
 (3.7,9.2)..controls(3.0,8.2)and(3.2,7.0)..(3.4,5.9)
 ..controls(3.5,4.9)and(3.8,4.0)..(4.2,4.0)
 ..controls(4.8,4.2)and(4.8,5.6)..(5.1,6.6)
 ..controls(5.3,7.7)and(4.9,8.8)..(4.4,9.2)--cycle;
\node[cgdepth] at (4.25,7.1) {85};
\node[cgdepth] at (2.45,6.0) {20};
\node[cgdepth] at (5.7,9.0) {20};
\node[cgdepth] at (4.8,3.2) {40};
\fill[white,draw=black,line width=.7pt]
 (5.0,8.6)..controls(5.2,8.9)and(5.4,8.5)..(5.35,8.1)
 ..controls(5.25,7.8)and(4.95,8.1)..(5.0,8.6)--cycle;
\node[cgnote] at (5.25,8.85) {island};
\cgspot{4.2}{12.2}{1}\cgspot{6.3}{10.55}{2}\cgspot{5.45}{9.25}{3}
\cgspot{4.85}{8.15}{4}\cgspot{1.75}{8.1}{5}\cgspot{1.65}{5.5}{6}
\cgspot{6.35}{4.65}{7}\cgspot{3.6}{3.0}{8}\cgspot{3.85}{1.2}{9}
\node[draw=black,fill=white,minimum size=3mm,inner sep=0pt,font=\scriptsize\bfseries]
 at (3.55,.35) {A};
\node[cgnote,anchor=north west] at (3.75,.45) {verify launch/signs};
\cgnorth{7.7}{12.6}\cgscale{.5}{.25}{1.55}{0.5 mile}
\end{tikzpicture}


\begin{minipage}{.93\linewidth}
\fontsize{7}{8.3}\selectfont
\textbf{Generalized depth contours:} solid 10 and 20 ft; long dash 40 ft;
short dash 60 ft; labeled maximum 85 ft. \textbf{Not for navigation.}
Independently drafted from the Wisconsin Conservation Department survey of
6 June 1955; present shoreline/access context checked against 2024 USDA/FSA
NAIP aerial orthophotography. No aerial image is reproduced. Numbered locations
are testable structure hypotheses, not promised fishing spots.
\end{minipage}
\end{center}

\newpage
\section{Nine places to begin the search}

\begin{longtable}{@{}p{.04\linewidth}p{.19\linewidth}p{.34\linewidth}p{.37\linewidth}@{}}
\toprule
\# & \textbf{Area} & \textbf{Why test it} & \textbf{First pass}\\
\midrule\endhead
1 & North narrows & Shallow water beside the first break can warm early. &
 Quietly fan-cast 4--12 ft; stop if no life is visible.\\
2 & Northeast points & Repeated hard-bottom points meet 10--20 ft water. &
 Cast across the break, then parallel it; note the exact depth of contact.\\
3 & North-central edge & A steep route between shelf and deep basin. &
 Hold the boat deep and work down the face; do not assume fish are on bottom.\\
4 & Island saddle & Shallow island structure adjoins basin water. &
 Circle once with electronics, then fish the windward edge and saddle.\\
5 & West shelf & The broadest gradual shelf on the historical sheet. &
 Find the outside vegetation line before choosing a rig.\\
6 & West pocket & Protected shallow-to-deep transition. &
 Check temperature, water clarity and new growth; leave if stagnant.\\
7 & Southeast turn & An inside turn interrupts the long eastern break. &
 Make depth-controlled passes from 10 to 25 ft.\\
8 & Southwest basin edge & Fast access between deep water and a tapering end. &
 Work the break at low light; scan open water for suspended bait.\\
9 & South narrows & Constricted shallow water near the recorded access. &
 Observe traffic and signs first; fish only where safe and lawful.\\
\bottomrule
\end{longtable}

\section{A disciplined Pine sequence}

\begin{enumerate}
\item Before launch, read the posted ordinance and access signs and check the
exact-water DNR regulation lookup.
\item Measure clarity and surface temperature. On bright, calm days begin
farther from the edge and make longer casts.
\item Search one shallow area, one break and one basin-adjacent feature. Record
depth, bottom, vegetation and bait before changing lure.
\item Treat historical cisco evidence as a reason to scan open water, not as
proof that a current cisco population or suspended predator bite exists.
\item Clean, drain and remove plants and animals before leaving; Pine has
multiple recorded aquatic invasive species.
\end{enumerate}

\begin{teloshazard}{Historical map, current water}
The source contours are more than seventy years old and generalized here.
Water level, vegetation, hazards, docks and access can change. Navigate with
current lawful aids and direct observation, never this sheet.
\end{teloshazard}

\newpage
\section{2026--27 rules printed for departure planning}

\begin{tabularx}{\linewidth}{@{}>{\bfseries}p{.22\linewidth}X@{}}
\toprule
Bass & Harvest May 2, 2026--Mar. 7, 2027; 14 in minimum; 5 daily.
Catch-and-release is open year-round unless otherwise noted.\\
Northern pike & May 2, 2026--Mar. 7, 2027; 26 in minimum; 2 daily.\\
Walleye group & May 2, 2026--Mar. 7, 2027; 18 in minimum; 3 daily.\\
Panfish & Open all year; no minimum; 25 daily.\\
Rock/yellow/white bass & Open all year; no minimum; unlimited.\\
Cisco/whitefish & Open all year; no minimum; 10 daily. Category does not prove presence.\\
Muskellunge & May 2--Dec. 31, 2026, open water only; 40 in minimum; 1 daily.
Category does not prove presence.\\
Motor trolling & Allowed with up to 3 hooks, baits or lures per angler.\\
\bottomrule
\end{tabularx}

\begin{teloshazard}{Recheck before every trip}
These are the Wisconsin DNR lookup results for WBIC 779200 checked 26 July
2026, not a permanent promise. Verify the exact water, species and date in the
live DNR rules. The DNR says regulation categories may not reflect fish actually
present. Posted landing signs are the definitive practical check for local
boating ordinances because the ordinance database may lag.
\end{teloshazard}

\section{Evidence boundary}

The current DNR lake record supports panfish, largemouth bass, smallmouth bass,
northern pike and walleye. Older agency-derived reports also record cisco,
common carp and two shiners, but historical presence is not current abundance.
The provider-neutral evidence record is
\texttt{research/lake-country-fishing/lakes-pine.md}.

\section{Primary sources}
\small
Wisconsin DNR Lake Page and contour sheet (WBIC 779200); Wisconsin DNR
exact-water fishing-regulation lookup (checked 2026-07-26); Wisconsin DNR
2026--27 regulation changes; USDA/FSA 2024 NAIP via the Wisconsin statewide
imagery service; SEWRPC Memorandum Reports 145 and 173. Full URLs and evidence
limits are recorded in the research file.

\end{document}
